%==============================================================================
% ProACT --- Proprioceptive Active Contact Sensing
% Target: IEEE-RAS Int. Conf. on Humanoid Robots (Humanoids) 2026
%   Format : IEEE conference (IEEEtran, two-column, US Letter, 10pt)
%   Limit  : 8 pages INCLUDING references
%   Deadline (paper): 2026-07-24 EOD Pacific ; optional video +3 days
% Build   : pdflatex -> bibtex -> pdflatex -> pdflatex
%           (needs IEEEtran.cls + IEEEtran.bst, both ship with TeX Live)
%==============================================================================
\documentclass[conference]{IEEEtran}
\IEEEoverridecommandlockouts

% --- packages ---------------------------------------------------------------
\usepackage{cite}
\usepackage{amsmath,amssymb,amsfonts}
\usepackage{algorithmic}
\usepackage{graphicx}
\usepackage{booktabs}      % nice tables (\toprule \midrule \bottomrule)
\usepackage{textcomp}
\usepackage{xcolor}
\usepackage{url}
\usepackage[hidelinks]{hyperref}

% small helper for TODO notes (remove before submission)
\newcommand{\todo}[1]{\textcolor{red}{[TODO: #1]}}

\def\BibTeX{{\rm B\kern-.05em{\sc i\kern-.025em b}\kern-.08em T\kern-.1667em\lower.7ex\hbox{E}\kern-.125emX}}

\begin{document}

\title{ProACT: Proprioceptive Active Contact Sensing for\\
       Whole-Body Contact Localization on Humanoids}

% ---- authors (placeholder) -------------------------------------------------
\author{
\IEEEauthorblockN{Anonymous Author(s)}
\IEEEauthorblockA{\textit{Affiliation} \\
\textit{Institution}\\
City, Country \\
email@domain}
% \and  ... more authors ...
}

\maketitle

%==============================================================================
\begin{abstract}
Proprioception-only estimation of external contact on humanoids has recently
reached whole-body wrench prediction, but its localization accuracy is
fundamentally limited: for contacts on the legs and torso during quiet stance
the external force lies largely in the null space of the contact Jacobian and
is therefore \emph{unobservable} from a fixed posture. We observe that contact
observability is a function of the robot's motion, and introduce \textbf{ProACT},
a proprioceptive \emph{active} contact-sensing framework: upon detecting a
\emph{sustained} contact whose location is ambiguous, the robot executes a small,
deliberate probing motion chosen to make the previously-unobservable force
component observable, and refines its estimate from the resulting proprioceptive
sequence. ProACT couples an always-on passive base estimator with an on-demand
active refinement stage, and is trained in simulation with ground-truth contact
labels.
% --- headline results (fill in) ---
\todo{one-sentence headline result: e.g. active probing raises leg/trunk
localization from X\% (passive) to Y\%, matching/exceeding the passive ceiling.}
\end{abstract}

\begin{IEEEkeywords}
humanoid robots, contact estimation, proprioception, active perception,
force sensing, whole-body control
\end{IEEEkeywords}

%==============================================================================
\section{Introduction}
% - the problem: sense whole-body external contact WITHOUT F/T or tactile sensors
% - why it matters: safe interaction, blind navigation, collision response
% - the gap: passive proprioceptive sensing is capped by motion-dependent
%   observability (legs/trunk); even large models plateau (cite SixthSense).
% - key idea: observability = f(motion); MOVE to make the unobservable observable.
% - requires SUSTAINED contact (transient contact is gone before you can probe).
% - contributions (bullet):
%     (1) an observability analysis of proprioceptive contact localization;
%     (2) ProACT: passive base + on-demand active probing (scripted -> learned);
%     (3) evidence that active probing breaks the passive ceiling on legs/trunk;
%     (4) a sim-to-real calibration path for deployment.
\todo{write introduction}

%==============================================================================
\section{Related Work}
\subsection{Proprioceptive external-force / contact estimation}
Model-based momentum observers recover external joint torques from dynamics
residuals~\cite{deluca2006momentum}. Learning-based methods now predict a
whole-body wrench field from proprioception alone~\cite{sixthsense}, and
self-supervised free-space torque models enable label-free calibration on real
robots~\cite{next}.
\todo{expand; position ProACT as the active complement to these passive methods}

\subsection{Active and tactile perception}
% active tactile exploration / information-gain probing (fingertip scale);
% dual control / active uncertainty reduction. Gap: none on whole-body
% humanoid proprioceptive contact.
\todo{}

\subsection{Whole-body compliance and interaction}
Impedance control~\cite{hogan1985impedance} underlies compliant physical
interaction; GentleHumanoid~\cite{gentlehumanoid} learns whole-body upper-body
compliance from proprioception without force/tactile sensors, but regulates a
manually-set stiffness rather than actively sensing where contact occurs.
\todo{expand}

%==============================================================================
\section{Problem Formulation and Observability}
% External force f at a body point produces joint torques tau_ext = J^T(q) f.
% Components of f in null(J^T(q)) produce zero torque -> unobservable in pose q.
% Stance masking (double support) + proximal ambiguity make legs/trunk hard.
% Key: null(J^T(q)) depends on q -> changing pose changes what is observable.
\begin{equation}
  \boldsymbol{\tau}_{\text{ext}} = \mathbf{J}^{\top}(\mathbf{q})\,\mathbf{f},
  \qquad
  \mathbf{f}\in\mathrm{null}\!\left(\mathbf{J}^{\top}(\mathbf{q})\right)
  \;\Rightarrow\; \boldsymbol{\tau}_{\text{ext}}=\mathbf{0}.
  \label{eq:jt}
\end{equation}
\todo{state the observability claim + why motion helps; optional Gramian.}

%==============================================================================
\section{Method}
\subsection{Passive base estimator}
% recurrent decoder over a proprioceptive window; outputs contact prob + per-link
% wrench FIELD (dense) rather than one-hot region. Detection-gated magnitude/dir.
\todo{}

\subsection{Active probing}
% trigger: sustained contact detected AND estimate ambiguous.
% probe = small command injected into a frozen low-level (impedance) policy:
% weight shift / torso sway / brief limb unload. Three tiers:
%   (a) scripted open-loop probe (existence proof)
%   (b) learned closed-loop probe policy (RL; reward = localization gain /
%       posterior-entropy reduction - balance/effort penalty)
%   (c) observability-Gramian / information-gain optimal probe (optional)
\todo{}

\subsection{Online refinement decoder}
% recurrent / recursive posterior over contact updated as the probe unfolds.
% prev_actions channel lets the model subtract self-motion from external signal.
\todo{}

\subsection{Sim-to-real calibration}
% NEXT-style free-space torque model on contact-free real data; bridges the
% probe policy trained in sim to the real robot.
\todo{}

%==============================================================================
\section{Experiments}
% Setup: humanoid in simulation, frozen stiff impedance low-level policy,
% net_pull applies SUSTAINED external force -> ground-truth (link + 3D force).
% Conditions: (P) passive static; (S) scripted probe; (A) learned probe.
% Controls: EQUAL total observation frames; same decoder capacity; same contact
% distribution. Metrics (comparable to prior work): exact-link acc, +/-1-link
% tolerant acc, detection rate, false-alarm rate; report legs/trunk separately.
% Mechanistic check: correlate localization gain with null-space excitation.
\todo{}

\begin{table}[t]
\centering
\caption{Contact localization: passive vs.\ active probing.}
\label{tab:main}
\begin{tabular}{lcccc}
\toprule
Method & \multicolumn{2}{c}{Exact-link (\%)} & \multicolumn{2}{c}{$\pm$1-link (\%)} \\
       & Legs/Trunk & Arms & Legs/Trunk & Arms \\
\midrule
Passive (static)      & \todo{} & \todo{} & \todo{} & \todo{} \\
ProACT (scripted)     & \todo{} & \todo{} & \todo{} & \todo{} \\
ProACT (learned)      & \todo{} & \todo{} & \todo{} & \todo{} \\
\bottomrule
\end{tabular}
\end{table}

%==============================================================================
\section{Results and Discussion}
\todo{}

%==============================================================================
\section{Conclusion}
\todo{}

%==============================================================================
% \section*{Acknowledgment}

\bibliographystyle{IEEEtran}
\bibliography{references}

\end{document}
